\documentclass[main.tex]{subfiles}

\begin{document}

\section{Endogenous Payment Choice}\label{sec:appendix-endogenous-choice}

When consumers choose payment methods, consumers also lose in aggregate from high interchange fees.
Intuitively, these losses reflect that too many consumers use high-fee methods to earn rewards.
The intuition behind the proof comes from taking revealed preference seriously.
Fix a consumer who is indifferent between debit cards and credit cards.
Consider a small increase in credit card rewards.
Such a consumer would then switch to credit cards to earn the additional rewards.
Because the consumer was indifferent to start with, she does not experience any increase in her private utility.
The greater rewards on the credit card exactly offset any non-pecuniary benefits (e.g. avoiding debt, lower adoption costs) of the debit card.
However, consumers collectively lose because the rewards are merely transfers funded by higher retail prices.

\subsection{Setup}\label{subsec:endo-setup}

We extend the consumer side of Appendix~\ref{sec:appendix-theory} to let consumers choose payment methods.
The game now has two stages.
In the first stage, consumers choose payment methods and merchants set retail prices; in the second stage consumers solve a consumption problem across merchants.
In contrast to the baseline model, there is now a continuum of consumers $i$ with measure one.
We describe the setup starting from the consumption problem.

\paragraph{Stage 2: Consumption.}
Consumer $i$ has utility function $U_{ik}$ over baskets of goods $q_{ijk}$ from merchants $j$.
Utility can depend on the payment method $k$, which captures the possibility that payment methods can have a causal effect on where consumers shop.
In this stage, consumer $i$ solves
\[
q^*_{ijk} = \arg\max_{q_{ijk}\geq 0} U_{ik}\left(\{q_{ijk}\}_{j}\right) \text{ s.t. } \sum_j p_j q_{ijk} \leq 1 + f_k,
\]
where $f_k$ is the reward for using payment method $k$.
The consumption problem induces an indirect utility function $\widetilde V_{ik}(p,f_k)$ and purchases $q^*_{ijk}(p,f_k)$.
Write $s^*_{ijk}\equiv p_j q^*_{ijk}(p,f_k)/(1+f_k)$ for her expenditure share at merchant $j$.

\paragraph{Stage 1: Consumer Payment Choice.}
Consumers have idiosyncratic tastes $\eta_i \sim F$ over payment methods.
We assume $F$ is atomless.
When choosing a payment method, consumer $i$ solves:
\[
k^*_i=\argmax_k\bigl\{\widetilde V_{ik}(p,f_k)+\eta_{ik}\bigr\}.
\]
Let $\pi_{ik}=\Pr(k^*_i=k\mid i)$ be the choice probability. 
To facilitate a comparison with the results with exogenous payment choice, we define the aggregate share of payment method $k$ to be $\mu_k=\int\pi_{ik}\,di$, and $\mu_{jk}=\frac{\int \pi_{ik}q^*_{ijk}\,di}{\sum_l \int \pi_{il}q^*_{ijl}\,di}$ to be method $k$'s share of merchant $j$'s revenue, so $\sum_k\mu_k=1$ and $\sum_k\mu_{jk}=1$.

\paragraph{Stage 1: Merchant pricing.}
Merchants price by~\eqref{eq:retail-price-pass} with $\rho^P=1$ and payment shares aggregated over all consumers; write $\hat\tau_j\equiv\sum_k\mu_{jk}\tau_{jk}$ for the payment-weighted average fee.\footnote{Because consumers and merchants are infinitesimal, the simultaneous move of prices and payment choices is without loss.}

\subsection{Theorem}\label{subsec:endo-theorem}

\begin{theorem}[Welfare effect under endogenous payment choice]\label{thm:endo-welfare}
Under Assumptions~\ref{assump:dominant-diagonal}--\ref{assump:income-elasticity} and atomless idiosyncratic tastes, the aggregate consumer welfare effect of a fee shift $dx_l$ is
\begin{equation}
\int\frac{dV_i}{\lambda_i}\,di = -\E_R[\sum_k \tau_{jk} d\mu_{jk}] + O(\tau_{\max}^2),
\label{eq:endo-total-welfare}
\end{equation}
each method's interchange fee times the change in its revenue share, averaged across merchants by revenue.
\end{theorem}

Within each merchant $\sum_k d\mu_{jk}=0$, so $\sum_k\tau_{jk}\,d\mu_{jk}=\sum_k(\tau_{jk}-c)\,d\mu_{jk}$ for any constant $c$: the loss is the covariance, across methods, between how expensive a method is and how much of the merchant's revenue it gains.
A reward-driven shift toward high-fee methods makes this positive and lowers welfare.

\subsection{Proof}\label{subsec:endo-proof}

\paragraph{Step A: Envelope theorem.}
By Roy's identity, conditional on choosing a payment method $k$ the consumer welfare effect is given by 
\begin{equation}
    \frac{dV_i}{\lambda_i}=-\sum_j s^*_{ijk^*_i}\,d\log p_j+df_{k^*_i}.
    \label{eq:endo-roy-perconsumer}
\end{equation}
By the Williams--Daly--Zachary discrete choice envelope theorem, this is also the change in ex-ante consumer welfare before choosing the payment method.
Integrating over consumers, the spending identity~\eqref{eq:spending-identity} gives $\int s^*_{ijk^*_i}\,di=\omega_j+O(\tau_{\max})$, and the reward term integrates to
\[
    \int df_{k^*_i}\,di = \sum_k \int df_k \, I\{k^*_i=k\}\,di=\sum_k\mu_k\,df_k.
\]
Together these give
\begin{equation}
\int\frac{dV_i}{\lambda_i}\,di=-\mathbb{E}_R[d\log p_j]+\sum_k\mu_k\,df_k+O(\tau_{\max}^2).
\label{eq:endo-roy-aggregate}
\end{equation}

\paragraph{Step B: Decompose and cancel.}
Applying the chain rule to the pricing formula~\eqref{eq:retail-price-pass} yields:
\begin{equation}
    d\log p_j = \sum_k \mu_{jk} d\tau_{jk} + \sum_k \tau_{jk} d\mu_{jk}
\end{equation}
By the full pass-through of interchange to rewards, we have that
\[
    df_k = \mu_k^{-1} \E_R[\mu_{jk} d\tau_{jk}] + O(\tau_{\max}^2)
\]
Substituting both expressions into \eqref{eq:endo-roy-aggregate} and dropping higher order terms yields:
\begin{equation}
    \int\frac{dV_i}{\lambda_i}\,di = -\E_R[\sum_k \tau_{jk} d\mu_{jk}] \qed
\end{equation}

\paragraph{Fees, not rewards.}
The loss is weighted by the interchange fee $\tau_{jk}$, not the reward $f_k$.
It is tempting to net the rewards back out---to argue that the interchange consumers pay returns to them as rewards---but it does not.
The rewards a \emph{switcher} earns are burned on the non-pecuniary utility she sacrifices, and the rewards an \emph{existing} user earns from a higher rate are spent paying the higher price that rate creates; the extra interchange from the reshuffling itself comes back to no consumer.
Of that interchange, the part funding rewards, $\E_R[\sum_k f_k\,d\mu_{jk}]$, is the non-pecuniary utility burned, and the remainder, $\E_R[\sum_k(\tau_{jk}-f_k)\,d\mu_{jk}]$, is the higher processing cost and issuer margin of the methods consumers move to. Consumers bear both.

\end{document}
